\section{Project Overview}

This report consolidates the four partial assignments into one antiship missile design path. The starting point is an Exocet/MM40-class equivalent missile constrained to solid propulsion. From that initial sizing, the model is carried through guidance, aerodynamic data generation, 6DOF preparation and autopilot design.

The work was developed in stages, with each stage passing data or design constraints to the next one:

\begin{table}[H]
\centering
\caption{Role of each partial report in the integrated project.}
\begin{tabular}{p{0.19\textwidth}p{0.32\textwidth}p{0.39\textwidth}}
\toprule
Stage & Main model & Role in the integrated project \\
\midrule
Report 1 & Preliminary sizing and 1D propulsion simulation & Defines mass, diameter, length, booster, sustainer, solid propellant and range capability. \\
Report 2 & Horizontal point-mass guidance simulation & First guidance version: pre-programmed trajectory, seeker logic and terminal switching. \\
Report 3 & DATCOM coefficients and 6DOF model & Converts geometry into aerodynamic data, CG/inertia and a higher-fidelity dynamic model. \\
Report 4 & Root-locus autopilot design & Uses the 6DOF/DATCOM data to design attitude and acceleration autopilots. \\
\bottomrule
\end{tabular}
\end{table}

In this sequence, Report 2 works as the first guidance prototype. Reports 3 and 4 then replace part of that simplified representation with DATCOM aerodynamics, mass properties and control design. The project flow is therefore:

\[
\text{Sizing} \rightarrow
\text{Point-mass guidance} \rightarrow
\text{DATCOM/6DOF dynamics} \rightarrow
\text{Autopilots}
\]

The resulting model has \(0.35~\text{m}\) diameter, \(5.8~\text{m}\) length, \(855~\text{kg}\) launch mass and simulated range above \(70~\text{km}\). It also includes horizontal-plane terminal guidance, a DATCOM aerodynamic database, a section-based CG/inertia model and autopilots satisfying the specified \(5g\) and \(0.5~\text{Hz}\) requirements.
